\paragraph{Integral generators after determinant extraction.}
\hypertarget{source-plane-symbols}{}
This part of the proof of Proposition~\ref{prop:residue-plane-neighbor-test}
gives the integral source polynomial on a single plane.
The determinant factors extracted above remain part of the
symbol. The result of the lifting is as follows.

\begin{lemma}[Lifting a plane symbol]\label{lem:lift-plane-symbol}
Fix a selected plane and the determinant factors of
Lemma~\ref{lem:plane-determinant-reserve}. Suppose that a residue symbol
\(\sigma\) has normal degree \(o\) and belongs to the prescribed
span of digit monomials times products of segment symbols,
with their prefix and multidegree conditions. Write \(\Pi\) for the selected plane. At precision \(c'\) there is one integral polynomial \(P\) with the following two restrictions
\[
 P|_\Pi\equiv w^o\sigma\pmod{w^{o+c'}},\qquad P|_{\Pi'}\equiv0\pmod{w^{o+c'}}\quad(\Pi'\text{ competing}).
\]
 It may be projected to the required
true multidegree. The column factors used in the lift have
prescribed order at most \(2Q\).
\end{lemma}

\begin{proof}
\proofstep{A paired path after the determinant extraction.}
Split the coordinate module into its plane and complement,
\[
 O^k=E\oplus C,\qquad E=O^r,\quad C=O^K,\quad K=k-r.
\]
At a top indexed by \(s\)
the even letters are vectors and the odd letters are covectors.
Fix one segment with end index \(z\). Its prescribed coarse
degree in a letter is the sum of its \(E\)-degree and its
\(C\)-degree. A segment symbol belongs to the span of the products \(A(E)G(C)\) satisfying the prefix bounds. To put the summands at the prescribed coarse
degrees, decompose \(A\) and \(G\) into homogeneous components
in every letter and keep the pairs whose degrees sum to the
prescribed ones. If \(\pi_\alpha\) extracts a multihomogeneous component, each prefix substitution \(g\) is linear within its letters. Hence
\[
 \pi_\alpha(gF)=g\,\pi_\alpha(F).
\]
Each component retains the same bounds. The
normal-degree cutoff is a condition on the original segment
symbol used in the rank argument. The realization here
applies to this whole span after the multidegree projection.

Apply the determinant extraction to \(A\) in each summand.
The determinants have fixed degrees in each coarse letter, so
the quotient \(A'\) also has fixed degrees. For every top \(s\)
after the incoming frame, let \(a_{<s},a_{\le s}\) denote the
minimum Laurent exponents of \(A'\) under the two prefix actions
ending before and at \(s\). Let \(g_{<s},g_{\le s}\) denote the
corresponding maximum exponents of \(G\). The remaining
separation says
\[
  g_{\le s}\le a_{<s}\quad(s\ {\rm even}),\qquad
  g_{<s}\le a_{\le s}\quad(s\ {\rm odd}).
\]
The complemented factor has degree zero in the incoming frame.

\proofstep{Integral basis after removing the determinant factors.}
Use the integral Pieri basis from
Section~\ref{sec:column-lifts} for the two factors separately.
The prefix filtration used there also works with a lower
bound on the last entry of the intermediate weight. In each
column the signed subset weight on one coordinate lies
between the contributions from the smallest and largest
entries of the intermediate weight. In the product it lies
between those two entries. This gives both bounds upon
conjugating the one-coordinate torus.

Conversely, over the
fraction field, the requirement that every conjugate torus
act within the specified range selects exactly the isotypic
summands whose first and last entries lie within that range.
Indeed the permitted subspace is invariant under the prefix
group. On each copy of \(V_\lambda\) in \(V_\lambda\otimes M_\lambda\), the extrema of the coordinate weights \(\mu\) are
\[
 \max\mu=\lambda_1,\qquad \min\mu=\begin{cases}\lambda_r,&\text{plane},\\\lambda_K,&\text{complement}.\end{cases}
\]
Complete reducibility gives the assertion. Splitting
successively by Pieri at the prefixes as in that section proves
the dimension formula for the intersection. The allowed basis
words span the intersection over \(O\) because the full basis is
unimodular.

Use the actual lower prefix exponents of \(A'\) as its lower
bounds and the actual upper prefix exponents of \(G\) as its
upper bounds, discarding zero summands. If \(\nu\) and
\(\kappa\) are the paths in any resulting product of basis
words, then for even \(s\) one has
\[
 \kappa_1(\le s)\le\nu_r(<s).
\]
For odd \(s\) one has
\(\kappa_1(<s)\le\nu_r(\le s)\).
These inequalities give exactly the interlacing across the
boundary between the \(r\)-th and \((r+1)\)-st rows of the
concatenated path \(\lambda=(\nu,\kappa)\). Each separate path already interlaces within its own rows.

For the boundary, if \(\lambda_i(s)\) denotes the combined
row weight after \(s\), the inequalities say
\[
 \lambda_r(s-1)\ge\lambda_{r+1}(s)\quad(s\ \mathrm{even}),
 \qquad
 \lambda_r(s)\ge\lambda_{r+1}(s-1)\quad(s\ \mathrm{odd}).
\]
The two row weights are nondecreasing at vector steps and
nonincreasing at covector steps. Thus the displayed boundary
inequality and the separate interlacings give dominance
both before and after each step. In the incoming frame
the complement stays at zero and plane entries are nonnegative
and increasing, so interlacing also holds there. Each strip's
size is the sum of the two sizes, and the concatenated path
starts at zero.

\paragraph{Full order of the two-block columns.}
\label{par:plane-column-orders}
Fix a column \(b\) of a concatenated path, and write \(h\) for
its height at a cut. A positive column uses the exterior
degree \(h\); a nonpositive column uses dual exterior degree
\(k-h\), as in Section~\ref{sec:column-lifts}. All changes of
height are by one, rising at a vector and falling at a
covector. For a subset \(I\) of size \(h\) in a positive column,
\[
 \#(I\cap C)=h-\#(I\cap E)\ge h-\min(r,h).
\]
Here the intersections count coordinate indices in the two summands. Equality characterizes the subsets of smallest normal cost. For a
nonpositive column
\[
 h>r
\]
at every cut. Indeed it starts at
\(k\) and loses at most one per letter, and the segment length
plus \(r\) is at most \(K\). The subsets of smallest cost
in the dual tensor use only complement indices.

The costly transitions have precisely the two forms
\[
 h\longrightarrow h+1\quad(h\ge r,\ \text{vector}),\qquad h\longrightarrow h-1\quad(h\le r,\ \text{covector}).
\]
In a nonpositive
column only the first possibility occurs. This is the normal
coordinate selected when both exterior subsets have their
prescribed plane occupancy. For a positive column, define the
defect of a subset as \(\min(r,h)\) minus its plane occupancy.
For a nonpositive column its defect is its plane occupancy.
In either case, along a sequence of wedge and contraction
terms, write \(n_i=1\) if the \(i\)-th letter
uses a normal coordinate and \(c_i=1\) if that step is
costly, taking both indicators to be zero in the other
case. The defect \(\delta_i\) then satisfies
\[
                    \delta_i-\delta_{i-1}=n_i-c_i .
\]
This follows from the change in occupancy in the wedge
or contraction operation and the change in the target
occupancy \(\min(r,h)\).

For example, at a positive rise from \(h<r\), adding a
complement index increases the defect by one. At a positive
fall from \(h\le r\), removing a complement index decreases
it by one. When \(h>r\) a rise on a plane index decreases the
defect and a fall on that index increases it. The dual
operations wedge on a covector and contract on a vector,
which increase and decrease its plane occupancy when the
plane index is used.

\proofstep{The allowed path of the restored column.}
The segment letters are the increasing enumeration of a set of
distinct top positions \(s_i\). Their radices
\[
 \omega_i=p^{z-s_i}>0
\]
are strictly decreasing. Let
\(\delta_i\) be the defect just after letter
\(i\), with \(\delta_0=0\). The excess of weighted normal cost
over the weighted count of costly steps is
\[
 \sum_{i<M}\delta_i(\omega_i-\omega_{i+1})+\delta_M\omega_M
\]
for a segment of \(M\) letters. Thus the lowest part is
computed by keeping defect zero throughout. In positive
columns the plane operations then take place with no
complement indices in the subset. All remaining operations
act on the complement past the plane indices.

In the
nonpositive columns the plane exterior degree in the dual
tensor stays zero. With the right wedge and contraction
conventions, an operation on a complement index counts only
larger complement indices for its sign, and every active
operation on a plane index counts only larger plane indices.
The defect-zero recursion therefore splits without a sign
change into the two column recursions for \(\nu\) and
\(\kappa\). For terminal coordinates take their union in the
corresponding blocks. Writing \(\operatorname{in}\) for the component of smallest normal degree, we obtain for the column polynomials \(C_b\)
\[
 \operatorname{in}\Bigl(\prod_b C_b\Bigr)=\prod_b\operatorname{in}(C_b).
\]
The defect-zero recursion identifies the right side with the prescribed paired word.

For the true lift, use the same exterior recursion in the
designated twisted frame. Between successive used indices
\(s<t\), raise every entry to \(p^{t-s}\), then perform the
operation at \(t\). After the last operation power up to \(z\).
Take \(\omega_i=p^{z-s_i}\), where \(s_i\) is the index of
letter \(i\). The column has degree \(\omega_i\) in each used
true letter. Its prescribed order is
\[
             d_b=\sum_{i\ {\rm costly}}\omega_i\le 2Q .
\]
The bound follows from the strict spacing of the used integer
indices and the geometric sum starting at the largest power,
which is at most \(Q\).

\proofstep{The order in a neighboring plane.}
For any tested plane, measure the running exterior tensor in
its own twisted frame. Suppose at a stage the common order is
at least \(d\), and every entry with positive defect has order
at least \(d+1\). This holds before the first operation.
For the passage \(s<t\) put \(P=p^{t-s}\) and write the
adapted coordinates at \(s\) as \(w^d U\). If \(L_s,L_t\)
are the changes of exterior coordinates from the designated
frame to the tested frame at the two times, then
\[
       (w^d U L_s^{-1})^{[P]} L_t
                 = w^{Pd}\big(U^{[P]}+pV\big)
\]
with \(V\) integral. Entrywise Frobenius is additive modulo
\(p\), and \((L_s^{-1})^{[P]}\equiv L_t^{-1}\pmod p\)
because the matrices and their inverse minors come from the
unramified ring and its successive Frobenius twists.

Since
\(p=w^2\) and \(P\ge2\), the terms in positive-defect entries
now have order at least \(Pd+2\). At a free step the free
coordinates take defect-zero subsets to defect-zero subsets,
and the normal coordinates contribute a factor \(w\).
At a costly step the free part kills the entries at defect
zero. The normal part carries them to the new defect-zero
subsets and contributes \(w\). Use as a common lower bound immediately before the linear step the order \(Pd\).
After a costly step the required order is \(d'=Pd+1\).
Thus the guaranteed lower bounds for the two contributions are
\[
 \begin{aligned}
 \text{defect zero, normal factor }w &: \quad Pd+1=d',\\
 \text{positive defect, integral factor} &: \quad Pd+2\ge d'+1 .
 \end{aligned}
\]
Only the first contribution can affect the leading coefficient.
At a free step the required order stays \(Pd\), and a normal
factor again raises it. These inequalities apply in every frame.

At the designated frame the change-of-frame matrices are
identities. After division by the prescribed power of \(w\)
the leading recursion simply keeps the defect-zero part. The
powering between steps and up to \(z\) commutes with its
additions and multiplications modulo \(w\), where the
characteristic is \(p\). Its result is therefore the
defect-zero coarse column with the substitution
\[
 X_{i,a}\mapsto X_{s_i,a}^{p^{z-s_i}}.
\]

Consequently the integral true column is a factor of order
\(d_b\) for the correction lemma. The product of its
normalized leading expressions is the desired paired symbol
after this substitution, with the full normal degree equal
to the sum of the column orders.  \label{cor:plane-symbols-1}

\paragraph{Digits, normal degree, and the final projection.}
Let \(X^\ell\) be the common digit monomial. At every top
each digit is less than the letter's radix, which is at
most \(Q\). Thus its total normal degree at that top is at
most \(kQ\). Spend one \(T_s\) for each normal digit factor
and use the leading order-one formulas for that coordinate.
If \(d_s\) normal digit factors are used at this frame, the reserve of \(T_s\) becomes
\[
 T_s^{kQ-d_s},\qquad kQ-d_s\ge0.
\]
It and the tangential digit variables are integral tangential coefficients. This
expresses \(X^\ell\prod_s T_s^{kQ}\) as a sum of products of
normalized leading expressions of universal order-one
factors with integral tangential coefficients. Every term has
the full normal degree of the digit monomial.

\proofstep{Recovering the true letter degree.}
Now start with a symbol \(\sigma\) of normal degree \(o\),
presented as \(X^\ell\) times a sum of products of the
transformed segment symbols. The extracted determinant
product is common to these summands. In the distinguished
coordinates it gives the identity
\[
  \sigma =
  F_*\,H\,X^\ell\!\prod_s T_s^{kQ}\,S + p^{c'}Y',
\]
where \(S\) is the sum of products of transformed residual
segments and \(Y'\) is integral. After expanding residual
segments in paired words, every summand on the right outside
the error has normal degree equal to the sum of the
appropriate column and digit orders. Since \(F_*,H\) are
tangential, take the component of normal degree \(o\).
The resulting main terms, after multiplication by \(w^o\),
are products of leading expressions of combined order \(o\),
with the reserves for precision \(c'\) contained in \(F_*\).
The error acquires order at least \(o+2c'\).

Apply the correction lemma to each surviving term, then add.
For the resulting integral true polynomial \(P\), the two tests read
\[
 \begin{aligned}P|_{\rm designated}&\equiv w^o\sigma\pmod{w^{o+c'}},\\ P|_{\rm competing}&\equiv0\pmod{w^{o+c'}}\qquad(\text{each competitor}).\end{aligned}
\]
If \(\sigma\) has a prescribed true multidegree,
project the resulting polynomial on that multidegree.
Every plane substitution is linear within each true
letter, so the projection commutes with all the substitutions
and preserves all divisibilities.  Its uniform order
bound on a key factor is \(2Q\), and its precision is the
prescribed \(c'\) at every plane. The determinant estimate above
also proves the one-sided version used for the dual test in
Section~\ref{par:plane-reassembly}. \label{cor:plane-symbols-2}
\end{proof}
